\documentclass{article}
\begin{document}
	First of, the main thing in any computer architecture, such as UCES, is its \textit{ISA}(Instruction Set Architecture).

	Translating the word \textit{ISA} to English: \textit{List of all commands your CPU can perform, for example, addition, substruction, memory operations, military grade encryption algorithms, etc}.

	UCES doesn't have some of listed above, but hopefully, it will be enough for most problems you want to solve. \newline

	At the next page you'll find all instruction in standard \textit{UCES}\newline

	But, now, I would like to do some FAQ:
	\begin{itemize}
		\item - Why do some instructions accept 3 arguments, while their analogs in \textit{x86\footnote{x86 is a reserved trademark by Intel}} accept only 2?\\
			- \textit{UCES} is not meant to be x86 compatible; Just read instruction definitions and everything will be clear :-)
		\item - You're referring to \textit{UCES} using words like \textit{architecture}, \textit{CPU}, does it mean \textit{UCES} has a physical implementation?\\
			- Nope. I haven't and won't make a physical silicon for \textit{UCES}, because I find it nonsensical, however, \textit{UCES} is open source and public domain, hence, you can freely design a sillicon for it if really want to.
	\end{itemize}

	\newpage
	\small
	\begin{enumerate}
		\item add $R_0$ $R_1$ $R_2$\\
			$R_2$ \textleftarrow \space $R_0$ + $R_1$
		\item sub $R_0$ $R_1$ $R_2$\\
			$R_2$ \textleftarrow \space $R_0$ - $R_1$
		\item and $R_0$ $R_1$ $R_2$\\
			$R_2$ \textleftarrow \space $R_0$ \& $R_1$
		\item or $R_0$ $R_1$ $R_2$\\
			$R_2$ \textleftarrow \space $R_0$ \textbar \space $R_1$
		\item xor $R_0$ $R_1$ $R_2$\\
			$R_2$ \textleftarrow \space $R_0$ \textasciicircum \space $R_1$
		\item nand $R_0$ $R_1$ $R_2$\\
			$R_2$ \textleftarrow \space \textasciitilde($R_0$ \& $R_1$)
		\item nor $R_0$ $R_1$ $R_2$\\
			$R_2$ \textleftarrow \space \textasciitilde($R_0$ \textbar \space $R_1$)
		\item xnor $R_0$ $R_1$ $R_2$\\
			$R_2$ \textleftarrow \space \textasciitilde($R_0$ \textasciicircum \space $R_1$)
		\item neg $R_0$ $R_1$\\
			$R_1$ \textleftarrow \space -$R_0$
		\item not $R_0$ $R_1$\\
			$R_1$ \textleftarrow \space \textasciitilde $R_0$
		\item div $R_0$ $R_1$ $R_2$\\
			$R_2$ \textleftarrow \space (u32)$R_0$/(u32)$R_1$
		\item idiv $R_0$ $R_1$ $R_2$\\
			$R_2$ \textleftarrow \space (s32)$R_0$/(s32)$R_1$
		\item rem $R_0$ $R_1$ $R_2$\\
			$R_2$ \textleftarrow \space (u32)$R_0$\%(u32)$R_1$
		\item irem $R_0$ $R_1$ $R_2$\\
			$R_2$ \textleftarrow \space (s32)$R_0$\%(s32)$R_1$
		\item mul $R_0$ $R_1$ $R_2$\\
			$R_2$ \textleftarrow \space (u32)$R_0$*(u32)$R_1$
		\item imul $R_0$ $R_1$ $R_2$\\
			$R_2$ \textleftarrow \space (s32)$R_0$*(s32)$R_1$
		\item wr\{8, 16, 32\} $R_0$ $R_1$\\
			$*R_0$ \textleftarrow \space $R_1$
		\item rd\{8, 16, 32\} $R_0$ $R_1$\\
			$R_1$ \textleftarrow \space *$R_0$
		\item rds\{8, 16\} $R_0$ $R_1$\\
			$R_1$ \textleftarrow \space (s32)*$R_0$
		\newpage
		\item cmp $R_0$ $R_1$ $R_2$\\
			$R_2$ \textleftarrow 0\\
			if $R_0$ = 0 $R_2$ \textleftarrow \space $R_2$ \textbar \space (1\textless\textless0);\\
			if $R_1$ = 0 $R_2$ \textleftarrow \space $R_2$ \textbar \space (1\textless\textless1);\\
			if $R_0$ = $R_1$ $R_2$ \textleftarrow \space $R_2$ \textbar \space (1\textless\textless2);\\
			if $R_0$ $\neq$ $R_1$ $R_2$ \textleftarrow \space $R_2$ \textbar \space (1\textless\textless3);\\
			if (s32)$R_0$ \textless \space (s32)$R_1$ $R_2$ \textleftarrow \space $R_2$ \textbar \space (1\textless\textless4);\\
			if (s32)$R_0$ \textgreater \space (s32)$R_1$ $R_2$ \textleftarrow \space $R_2$ \textbar \space (1\textless\textless5);\\
			if $R_0$ \textless \space $R_1$ $R_2$ \textleftarrow \space $R_2$ \textbar \space (1\textless\textless6);\\
			if $R_0$ \textgreater \space $R_1$ $R_2$ \textleftarrow \space $R_2$ \textbar \space (1\textless\textless7);\\
			if carry($R_0$ - $R_1$) $R_2$ \textleftarrow \space $R_2$ \textbar \space (1\textless\textless8);\\
			if carry($R_0$ + $R_1$) $R_2$ \textleftarrow \space $R_2$ \textbar \space (1\textless\textless9);\\
			if overflow($R_0$ * $R_1$) $R_2$ \textleftarrow \space $R_2$ \textbar \space (1\textless\textless10)
		\item mv $R_0$ $R_1$\\
			$R_1$ \textleftarrow $R_0$
		\item mve $R_0$ $R_1$ $R_2$\\
			if $R_0$ $\neq$ 0 $R_2$ \textleftarrow $R_1$
		\item mvo $R_0$ $R_1$ $R_2$\\
			if $R_0$ = 0 $R_2$ \textleftarrow $R_1$
		\item tsbi $R_0$ $I_1$ $R_2$\\
			if $R_0$ \& (1\textless\textless$I_1$) = 0 $R_2$ \textleftarrow \space $R_2$ \textbar \space 0;\\
			else $R_2$ \textleftarrow \space $R_2$ \textbar \space 1
		\item tsi $R_0$ $I_1$ $R_2$\\
			if $R_0$ \& $I_1$ = 0 $R_2$ \textleftarrow \space $R_2$ \textbar \space 0;\\
			else $R_2$ \textleftarrow \space $R_2$ \textbar \space 1
		\item tsb $R_0$ $R_1$ $R_2$\\
			if $R_0$ \& (1\textless\textless$R_1$) = 0 $R_2$ \textleftarrow \space $R_2$ \textbar \space 0;\\
			else $R_2$ \textleftarrow \space $R_2$ \textbar \space 1
		\item ts $R_0$ $R_1$ $R_2$\\
			if $R_0$ \& $R_1$ = 0 $R_2$ \textleftarrow \space $R_2$ \textbar \space 0;\\
			else $R_2$ \textleftarrow \space $R_2$ \textbar \space 1
		\item call $R_0$\\
			push(pc+4);\\
			pc = $R_0$
		\item ret\\
			pc = pop()
		\item push $R_0$\\
			push($R_0$)
		\item pop $R_0$\\
			$R_0$ = pop()
		\item lsl $R_0$ $R_1$ $R_2$\\
			$R_2$ \textleftarrow \space $R_0$\textless\textless$R_1$
		\item lsr $R_0$ $R_1$ $R_2$\\
			$R_2$ \textleftarrow \space $R_0$\textgreater\textgreater$R_1$
		\item andi $R_0$ $I_1$ $R_2$\\
			$R_2$ \textleftarrow \space $R_0$\&$I_1$\&0xFF
		\item ori $R_0$ $I_1$ $R_2$\\
			$R_2$ \textleftarrow \space $R_0$\textbar$I_1$\&0xFF
		\item debug ???\\
			something()
			// Effect is decided purely by implementation. Don't use this instruction in production code
		\item addi $R_0$ $I_1$ $R_2$\\
			$R_2$ \textleftarrow \space $R_0$+(s8)$I_1$
		\item hlt\\
			halt() // Stops execution of current program.
		\item udi\\
			trap(undefined instruction)
		\item ldlx $I_0$ $R_1$\\
			$R_1$ \textleftarrow \space $I_0$\&0xFFFF
		\item ldhx $I_0$ $R_1$\\
			$R_1$ \textleftarrow \space $I_0$\&0xFFFF \textless\textless 16
		\item ldl $I_0$ $R_1$\\
			$R_1$ \textleftarrow \space $R_1$\&0xFFFF0000 \textbar \space $I_0$\&0xFFFF
		\item ldh $I_0$ $R_1$\\
			$R_1$ \textleftarrow \space $R_1$\&0xFFFF \textbar \space ($I_0$\&0xFFFF \textless\textless 16)
	\end{enumerate}

\end{document}
